% Part: sets-functions-relations
% Chapter: induction
% Section: induction

\documentclass[../../../include/open-logic-section]{subfiles}

\begin{document}

\olfileid{sfr}{ind}{int}

\olsection{د طبيعي عددونو او فورمولونو استقرا}

\begin{explain}
استقرا د رياضياتو د ثبوت يوه عامه طريقه ده چې په منطق کښې هم کارېږي۔ ښايي له رياضياتو سره ئې اشنا يئ، چې پر طبيعي عددونو استقرا کوئ۔ د رياضيکي او منطقي څيزونو ډېر نور ډولونه هم داسې دي چې استقرا پرې کېدے شي۔

په عمومي ډول په استقرا يوه کلي دعوه ثابتوو؛ يعنې ښيو چې د يو ټاکلي ډول هر څيز يو خاصيت لري۔ استقرا په تېره بيا هغه وخت ګټوره وي چې د کلي حکم د مستقيم داخلولو طريقه ډېره پېچلې وي۔ د جوړښتي استقرا دپاره اړوند څيزونه په ځانګړي ډول جوړېږي: که د سټ هر غړے يا بنسټيز وي يا په متناهي ګامونو له بنسټيزو غړو جوړ شوے وي، پر دغه جوړښت استقرا کولے شو۔ لاندې تعريفونه دا خبره څرګندوي۔
\end{explain}

\begin{defn}
سټ $S$ د تابعې $f$ لاندې هله \emph{تړلے} دے چې تابع ئې پر ټولو غړو تعريف وي او، هر کله چې $x \in
S$، نو $f(x) \in S$۔
\end{defn}

\begin{explain}
د بېلګې په توګه، طبيعي عددونه ($\mathbb{N}$) د ځاے‌ناستې تابعې $f(x) = x+1$ لاندې تړلي دي۔ که $x$ طبيعي عدد وي، $x+1$ هم طبيعي عدد دے۔ خو د مخکيني عدد د تابعې $f(x) = x-1$ لاندې تړلي نۀ دي، ځکه $0$ طبيعي عدد دے خو $-1$ طبيعي عدد نۀ دے۔
\end{explain}

\begin{defn}
وايو چې خاصيت $P$ د تابعې $f(x)$ لاندې \emph{ساتل کېږي} که د $f$ د تعريف په ساحه کښې د هر څيز $a$ دپاره له $P(a)$ څخه $P(f(a))$ راځي؛ يعنې که $a$، $P$ خاصيت ولري، $f(a)$ ئې هم لري۔
\end{defn}

\begin{explain}
د «جفت عدد کېدو» خاصيت د تابعې $f(x) = x+2$ لاندې ساتل کېږي، ځکه که $x$ جفت وي، $x+2$ هم جفت دے۔ خو د ځاے‌ناستې تابعې لاندې نۀ ساتل کېږي؛ که $x$ جفت وي، $x+1$ طاق دے۔

پر طبيعي عددونو استقرا ځکه کېږي چې له صفر نه د ځاے‌ناستې تابعې په متناهي ځلو کارولو هر طبيعي عدد ته رسېږو۔ په همدې ډول د جوړښتي استقرا وړ سټونو غړي هم له خپل بنسټ نه په متناهي ګامونو جوړېږي۔
\end{explain}

\begin{defn}[پر $\Nat$ استقرا]
که صفر $P$ خاصيت ولري او $P$ د ځاے‌ناستې تابعې لاندې وساتل شي، نو هر طبيعي عدد $P$ خاصيت لري۔
\end{defn}

\begin{ex} د لومړنيو $n$ مثبتو طبيعي عددونو جمع $\frac{n(n+1)}{2}$ ده۔ \\

\textbf{بنسټيز حالت۔} د صفر عددونو جمع $0=
\frac{0\cdot 1}{2}$ ده۔\\

\textbf{استقرايي ګام۔} فرض کړئ چې خاصيت د $k$ دپاره سم دے۔ ښيو چې د $k+1$ دپاره هم سم دے۔ يعنې فرض کړئ $P(k)$ صدق کوي:
\[ 0 + 1 + \ldots + k = \frac{k(k+1)}{2} \]
ښيو چې $P(k+1)$ صدق کوي، يعنې:
\[ 0 + 1 + \ldots + k + (k+1) = \frac{(k+1)(k+2)}{2} \]
د دې دپاره هغه فرض کاروو چې $P(k)$۔

\begin{align*}
0 + 1 + \ldots + k &= \frac{k(k+1)}{2} \tag{فرض} \\
0 + 1 + \ldots + k + (k+1) &= \frac{k(k+1)}{2} + (k+1)
\tag{دواړو خواوو ته $k+1$ زيات کړئ} \\
&= \frac{k(k+1) + 2(k+1)}{2} \\
&= \frac{k^2 + k + 2k +2}{2} \\
&=\frac{(k+1)(k+2)}{2}
\end{align*}
نو استقرايي ګام پوره شو او دعوه ثابته شوه۔
\end{ex}

\begin{explain}
د !!{formula}s بنسټيز غړي اټومي !!{formula}s دي۔ پېچلي !!{formula}s د لږ پېچلو !!{formula}s نه د منطقي عملګرو په کارولو جوړېږي۔ هر عملګر له داسې تابعې سره سمون لري چې له اړينو !!{formula}s څخه يو نوے !!{formula} جوړوي او هماغه عملګر پرې کاروي۔ د بېلګې په توګه، هغه تابع چې !!{formula}s $!A$ او $!B$ د او په تړنه تړي، $\mathcal E _\land (!A,
!B) = !A \land !B$ ليکلے شو۔ دې تابعو ته \emph{فورمول جوړوونکې تابعې} وايو۔ د !!{formula}s سټ د دغو تابعو لاندې تړلے دے؛ هر کله چې $!A$ او $!B$ !!{formula}s وي، $\lnot !A$، $!A \land
!B$ او داسې نور هم فورمولونه دي۔
\end{explain}

\begin{defn}[پر !!{formula}s استقرا]
که هر اټومي !!{formula}، $P$ خاصيت ولري او $P$ د !!{formula} جوړوونکو تابعو لاندې وساتل شي، نو هر !!{formula}، $P$ خاصيت لري۔
\end{defn}

\begin{explain}
دا تعريف پر !!{formula}s د استقرايي ثبوت يوه طريقه راښيي۔ که ښودل غواړو چې هر !!{formula}، $P$ خاصيت لري، لاندې ګامونه اخلو۔\\

\textbf{بنسټيز حالت۔} فرض کړئ $!A$ يو اټومي !!{formula} دے۔ [\ldots]
نو $!A$، $P$ خاصيت لري۔

\textbf{استقرايي ګام۔} فرض کړئ $!A$ او $!B$ داسې !!{formula}s دي چې دواړه $P$ خاصيت لري۔

د $\lnot$ حالت۔ [\ldots] نو $\lnot !A$، $P$ خاصيت لري۔

د $\land$ حالت۔ [\ldots] نو $!A \land !B$، $P$ خاصيت لري۔

د $\lor$ حالت۔ [\ldots] نو $!A \lor !B$، $P$ خاصيت لري۔

د $\lif$ حالت۔ [\ldots] نو $!A \lif !B$، $P$ خاصيت لري۔

د $\liff$ حالت۔ [\ldots] نو $!A \liff !B$، $P$ خاصيت لري۔

د $\lforall[]$ حالت۔ [\ldots] نو $\lforall[x] !A$، $P$ خاصيت لري۔

د $\lexists[]$ حالت۔ [\ldots] نو $\lexists[x] !A$، $P$ خاصيت لري۔ \\

نو هر !!{formula}، $P$ خاصيت لري۔
\end{explain}

\end{document}
